\documentclass[11pt,a4paper]{article}
\usepackage[UTF8]{ctex}
\usepackage{geometry}
\geometry{left=2.5cm,right=2.5cm,top=2.5cm,bottom=2.5cm}
\usepackage{graphicx}
\usepackage{amsmath}
\usepackage{booktabs}
\usepackage{hyperref}
\usepackage{caption}
\usepackage{enumitem}
\usepackage{fancyhdr}
\usepackage{titlesec}

\pagestyle{fancy}
\setlength{\headheight}{14pt}  
\fancyhf{}
\fancyhead[C]{实验报告：Cache优化与指令级并行对程序性能的影响}
\fancyfoot[C]{\thepage}

\titleformat{\section}{\large\bfseries}{\thesection}{1em}{}
\titleformat{\subsection}{\normalsize\bfseries}{\thesubsection}{1em}{}

\hypersetup{
    colorlinks=true,
    linkcolor=blue,
    urlcolor=blue,
}

\title{\textbf{实验报告：Cache优化与指令级并行对程序性能的影响}}
\author{邓耀飞 \quad 学号：2413404}
\date{}

\begin{document}

\maketitle
\thispagestyle{empty}

\begin{abstract}
本实验旨在探究现代计算机体系结构中Cache局部性和指令级并行对程序性能的影响。通过实现矩阵列内积的两种算法（逐列访问 vs 按行访问）和数组求和的两种算法（链式累加 vs 两路累加），在不同数据规模下进行性能测试。实验结果表明，Cache优化算法在矩阵规模为4096时获得5.68倍加速比，两路累加在小规模数组中获得约2倍加速比。本文还讨论了编译器优化对性能测量的影响，并提出了正确的测量方法。源代码已上传至GitHub仓库：\url{https://github.com/K1Ta-Ynni/hpc-performance-lab}
\end{abstract}

\tableofcontents
\newpage

\section{实验目的}

\begin{enumerate}
    \item 理解Cache局部性对矩阵运算性能的影响，对比“逐列访问”与“按行访问”两种算法的执行效率。
    \item 理解指令级并行对数组求和性能的影响，对比“链式累加”与“两路累加”两种算法的执行效率。
    \item 掌握在编译器优化（\texttt{-O2}）下进行性能测试的正确方法，避免“死代码删除”带来的测量误差。
\end{enumerate}

\section{实验平台}

\begin{table}[h]
\centering
\caption{实验平台配置}
\begin{tabular}{|l|l|}
\hline
\textbf{项目} & \textbf{配置} \\
\hline
CPU & Intel Core i7-14650HX @ 2.20GHz \\
\hline
内存 & 16.0 GB \\
\hline
显卡 & NVIDIA GeForce RTX 4060 Laptop GPU \\
\hline
操作系统 & Windows 11 家庭中文版 (24H2) \\
\hline
编译器 & g++ 8.1.0 (MinGW-W64) \\
\hline
优化级别 & \texttt{-O2} \\
\hline
计时方式 & \texttt{std::chrono::high\_resolution\_clock}，自适应重复次数 \\
\hline
\end{tabular}
\end{table}

\section{实验内容与算法设计}

\subsection{矩阵列内积}

\textbf{问题描述}：计算 \(n \times n\) 矩阵 \(A\) 与向量 \(b\) 的每一列内积：
\[
\text{result}[j] = \sum_{i=0}^{n-1} A[i][j] \times b[i]
\]

\textbf{算法设计}：

\begin{itemize}
    \item \textbf{平凡算法（逐列访问）}：外层循环遍历列，内层循环遍历行。内存访问模式为 \(A[0][j], A[1][j], A[2][j]\ldots\)，空间局部性差。
    \item \textbf{Cache优化算法（按行访问）}：外层循环遍历行，内层循环遍历列，利用局部变量累加每列结果。内存访问模式为 \(A[i][0], A[i][1], A[i][2]\ldots\)，空间局部性好，充分利用Cache line。
\end{itemize}

\subsection{数组求和}

\textbf{问题描述}：计算长度为 \(n\) 的数组所有元素之和。

\textbf{算法设计}：

\begin{itemize}
    \item \textbf{平凡算法（链式累加）}：单变量依次累加，每次迭代依赖上一次结果，无法利用指令级并行。
    \item \textbf{两路累加（指令级并行优化）}：使用两个累加变量，分别处理偶数位和奇数位元素，最后合并。减少数据依赖，允许CPU同时执行两条加法指令。
\end{itemize}

\section{实验方案}

\subsection{测试规模}

\begin{table}[h]
\centering
\caption{测试规模}
\begin{tabular}{|l|l|}
\hline
\textbf{测试项} & \textbf{规模范围} \\
\hline
矩阵列内积 & n = 256, 512, 1024, 2048, 4096 \\
\hline
数组求和 & n = 100000, 500000, 1000000, 5000000, 10000000 \\
\hline
\end{tabular}
\end{table}

\subsection{计时方法}

采用自适应重复次数策略：
\begin{enumerate}
    \item 先估算单次执行时间
    \item 自动调整重复次数，使总测量时间稳定在约200-500ms
    \item 取单次平均时间作为最终结果
\end{enumerate}

\subsection{防止编译器优化}

在 \texttt{-O2} 优化级别下，未被使用的返回值会被编译器视为“无用代码”并删除。解决方案：将计算结果赋值给 \texttt{volatile} 变量，强制编译器保留计算过程。

\section{实验结果与分析}

\subsection{矩阵列内积性能对比}

\begin{table}[h]
\centering
\caption{矩阵列内积性能对比}
\begin{tabular}{|c|c|c|c|}
\hline
\textbf{n} & \textbf{平凡算法 (ms)} & \textbf{Cache优化 (ms)} & \textbf{加速比} \\
\hline
256 & 0.041 & 0.024 & 1.72 \\
\hline
512 & 0.204 & 0.091 & 2.24 \\
\hline
1024 & 1.030 & 0.359 & 2.87 \\
\hline
2048 & 9.484 & 2.371 & 4.00 \\
\hline
4096 & 58.169 & 10.242 & \textbf{5.68} \\
\hline
\end{tabular}
\end{table}

\textbf{分析}：
\begin{enumerate}
    \item 加速比随 \(n\) 增大而显著提升，从 1.72x 增长至 5.68x。
    \item 当 \(n=4096\) 时，矩阵大小约为 \(4096^2 \times 8\ \text{B} = 128\ \text{MB}\)，远超 L3 Cache 容量（通常 8–32 MB）。平凡算法因频繁 Cache Miss 性能严重下降，而 Cache 优化算法仍保持较高效率。
    \item 平凡算法按列访问时，相邻元素在内存中相隔一个列宽（约 32KB），破坏空间局部性；按行访问则连续读取，每次 Cache Line（64 字节）加载的 8 个元素全部被使用。
\end{enumerate}

\subsection{数组求和性能对比}

\begin{table}[h]
\centering
\caption{数组求和性能对比}
\begin{tabular}{|c|c|c|c|}
\hline
\textbf{n} & \textbf{平凡算法 (μs)} & \textbf{两路累加 (μs)} & \textbf{加速比} \\
\hline
100000 & 43.85 & 22.59 & 1.94 \\
\hline
500000 & 217.70 & 109.33 & 1.99 \\
\hline
1000000 & 437.80 & 222.65 & 1.97 \\
\hline
5000000 & 3050.78 & 2629.95 & 1.16 \\
\hline
10000000 & 6622.43 & 4900.88 & 1.35 \\
\hline
\end{tabular}
\end{table}

\textbf{分析}：
\begin{enumerate}
    \item 在 \(n \leq 10^6\) 时，两路累加稳定达到约 \textbf{2 倍加速比}。这是因为两路累加消除了数据依赖，允许 CPU 的两个浮点运算单元同时工作，达到指令级并行的理论极限。
    \item 当 \(n\) 增大至 \(5 \times 10^6\) 以上时，加速比下降至 1.16–1.35。此时内存访问时间成为主要瓶颈，计算优势被掩盖。
    \item 验证输出显示两种算法结果一致（均为 \(n\)），证明算法实现正确。
\end{enumerate}

\section{进阶实验与讨论}

\subsection{编译器优化对性能测量的影响}

在 \texttt{-O2} 级别下，未使用的函数返回值会被编译器视为“无用代码”并删除。实验中，数组求和函数因仅返回 \texttt{double} 且未使用，被完全优化，测量结果为 0。

\textbf{解决方法}：使用 \texttt{volatile} 关键字强制保留计算结果：
\begin{verbatim}
volatile double result = Sum::naive(arr);
\end{verbatim}

这体现了在高性能计算中，开发者必须理解编译器的优化行为，才能正确评估算法性能。

\subsection{性能瓶颈分析}

\begin{table}[!htb]
\centering
\caption{性能瓶颈分析}
\label{tab:bottleneck}
\begin{tabular}{|l|l|l|}
\hline
\textbf{问题} & \textbf{瓶颈类型} & \textbf{优化效果} \\
\hline
矩阵列内积 & Cache Miss & 5.68x 加速比 \\
\hline
数组求和（小规模） & 指令级并行 & 约 2x 加速比 \\
\hline
数组求和（大规模） & 内存带宽 & 加速比下降至 1.3x \\
\hline
\end{tabular}
\end{table}

从表\ref{tab:bottleneck}可以看出，不同计算任务在不同规模下的性能瓶颈存在显著差异：

\textbf{矩阵列内积（Cache Miss 主导）}：当矩阵规模为 4096 时，数据量达到 128MB，远超 CPU 的 L3 Cache 容量（通常 8-32MB）。平凡算法因逐列访问导致频繁的 Cache Miss，每次访问都可能触发从内存加载数据，耗时显著增加。而 Cache 优化算法通过按行访问，充分利用空间局部性，在 Cache 命中率上获得巨大优势，从而取得 5.68 倍加速比。

\textbf{数组求和（小规模：指令级并行有效）}：当数组规模不超过 100 万时，数据可以完全放入 Cache，内存访问时间不是瓶颈。此时，两路累加通过减少数据依赖，允许 CPU 的两个浮点运算单元同时工作，实现了接近 2 倍的理论加速比。这验证了现代超标量处理器的指令级并行能力。

\textbf{数组求和（大规模：内存带宽瓶颈）}：当数组规模达到 500 万以上时，数据总量超出 Cache 容量，每次求和都需要从内存读取数据。此时内存带宽成为主要限制因素，计算速度被数据供给速度拖慢。因此两路累加的计算优势被掩盖，加速比下降至 1.16-1.35 倍。

上述分析表明，性能优化策略需要根据问题规模和硬件特性灵活选择：计算密集型小规模问题适合指令级并行优化，而数据密集型大规模问题则需要优先考虑 Cache 优化和内存访问模式优化。

\section{实验结论}

\begin{enumerate}
    \item \textbf{Cache 优化效果显著}：按行访问相比按列访问，在大规模矩阵运算中最高可获得 \textbf{5.68 倍} 加速比。性能提升来源于更好的空间局部性和更高的 Cache 命中率。
    \item \textbf{指令级并行有效}：两路累加在小规模数组中实现约 \textbf{2 倍} 加速比，验证了超标量处理器并行执行多条独立指令的能力。
    \item \textbf{性能瓶颈随规模变化}：小规模时计算能力是瓶颈，大规模时内存带宽成为主要限制因素。
    \item \textbf{正确测量需考虑编译器优化}：在高优化级别下，需通过 \texttt{volatile} 等手段防止被测代码被优化掉，确保测量结果真实有效。
\end{enumerate}

\section{源代码}

本次实验的完整源代码已上传至 GitHub 仓库：
\begin{center}
\url{https://github.com/K1Ta-Ynni/hpc-performance-lab}
\end{center}

仓库包含：
\begin{itemize}
    \item \texttt{benchmark.cpp} —— 完整的性能测试代码
    \item \texttt{report.tex} —— 本实验报告源文件
    \item \texttt{README.md} —— 项目说明文档
\end{itemize}

\section*{参考文献}

\begin{enumerate}[label={[1]}]
    \item Hennessy, J. L., \& Patterson, D. A. (2017). \textit{Computer Architecture: A Quantitative Approach} (6th ed.). Morgan Kaufmann.
    \item 陈国良, 吴俊敏. (2018). \textit{并行计算机体系结构}. 高等教育出版社.
    \item Intel Corporation. (2024). \textit{Intel Core i7-14650HX Processor Specifications}.
    \item GCC Team. (2023). \textit{GCC Optimization Options}. \url{https://gcc.gnu.org/onlinedocs/gcc/Optimize-Options.html}
\end{enumerate}

\end{document}